\subsubsection*{定理 3.4.}
设 $X$ 是 $\mathbb F_q$ 上的有限型分离概形.
以 $H_c^i$ 表示带固有支撑的平展上同调 (SGA~4, XVII~5.1.9(iii)), 则有

\begin{equation}\tag{3.4.1}
Z(t,X/\mathbb F_q)
\equiv
\prod_{i\geq0}
\det\!\left(1-t\varphi_q^{-1}\mid
H_c^i(\bar X,\mathbb Z/p\mathbb Z)\right)^{(-1)^{i+1}}
\pmod p.
\end{equation}

先证明 (3.4.1) 与 (3.1.1) 等价. 暂记

\begin{equation}\tag{3.4.2}
Z(X)\quad\text{替代 }Z(t,X/\mathbb F_q)
\end{equation}

及

\begin{equation}\tag{3.4.3}
Z'(X)\quad\text{替代 }\quad
\prod_i\det\!\left(1-t\varphi_q^{-1}\mid
H_c^i(\bar X,\mathbb Z/p\mathbb Z)\right)^{(-1)^{i+1}}.
\end{equation}

\subsubsection*{引理 3.5.}
若 $\{U_i\}$ 是由开子概形组成的 $X$ 的有限覆盖, 则

\begin{equation}\tag{3.5.1}
Z'(X)=
\prod_{j\geq0}
\prod_{1\leq i_0<\cdots<i_j\leq n}
Z'(U_{i_0}\cap\cdots\cap U_{i_j})^{(-1)^j}.
\end{equation}

\medskip
\noindent\emph{证明.}
这立即由局部化谱序列 (SGA~4, XVII~6.2.10) 得出

\[
E_1^{-a,b}=\bigoplus_{i_0<\cdots<i_a}
H_c^b(\bar U_{i_0}\cap\cdots\cap\bar U_{i_a},\mathbb Z/p\mathbb Z)
\Longrightarrow H_c^{-a+b}(\bar X,\mathbb Z/p\mathbb Z).
\]

\begin{editorialnote}
原书第412页及英文底本在此将交集的末指标写成 $i_b$, 上同调次数写成 $a$.
在保留 $E_1^{-a,b}$ 及总次数 $-a+b$ 的约定下, 交集应有 $a+1$ 项, 上同调次数应为 $b$;
上式据此交换这两处指标.
\end{editorialnote}

亦可利用下面的引理~3.6 对 $n$ 归纳得到.

\subsubsection*{引理 3.6.}
若 $U$ 为 $X$ 的开子概形, $Y$ 为闭子概形 $X-U$, 则

\begin{equation}\tag{3.6.1}
Z'(X)=Z'(U)Z'(Y).
\end{equation}

\medskip
\noindent\emph{证明.}
这由上同调正合列得出

\begin{equation}\tag{3.6.2}
\longrightarrow H_c^i(\bar U,\mathbb Z/p\mathbb Z)
\longrightarrow H_c^i(\bar X,\mathbb Z/p\mathbb Z)
\longrightarrow H_c^i(\bar Y,\mathbb Z/p\mathbb Z)\longrightarrow.
\end{equation}

\subsubsection*{引理 3.7 (M\"obius).}
设 $B$ 为有限集, $\{A_i\}$ 为由子集组成的有限覆盖. 则

\begin{equation}\tag{3.7.1}
\operatorname{card}(B)=
\sum_{j\geq0}(-1)^j
\sum_{i_0<\cdots<i_j}
\operatorname{card}(A_{i_0}\cap\cdots\cap A_{i_j}).
\end{equation}

\medskip
\noindent\emph{证明.}
对于离散空间 $B$ 的有理上同调和开覆盖 $\{A_i\}$, 有 Leray 谱序列

\[
E_1^{p,q}=\bigoplus_{i_0<\cdots<i_q}
H^p(A_{i_0}\cap\cdots\cap A_{i_q},\mathbb Q)
\Longrightarrow H^{p+q}(B,\mathbb Q),
\]

随后计算 $B$ 的 Euler--Poincar\'e 示性数即可.

\subsubsection*{引理 3.8.}
若 $\{U_i\}$ 为由子概形组成的 $X$ 的有限覆盖, 则

\begin{equation}\tag{3.8.1}
Z(X)=
\prod_{j\geq0}
\prod_{i_0<\cdots<i_j}
Z(U_{i_0}\cap\cdots\cap U_{i_j})^{(-1)^j}.
\end{equation}

\medskip
\noindent\emph{证明.}
有

\begin{equation}\tag{3.8.2}
\operatorname{card}(X(\mathbb F_{q^n}))
=\sum_{r\mid n}r\,\operatorname{card}\{y\in X\mid\deg(y)=r\},
\end{equation}

这给出

\begin{equation}\tag{3.8.3}
Z(X)=\exp\!\left(
\sum_{n\geq1}\operatorname{card}(X(\mathbb F_{q^n}))\frac{t^n}{n}
\right).
\end{equation}

结论由 (3.7.1) 得出.

\subsubsection*{引理 3.9.}
假设公式 (3.4.1), 或等价的 (3.1.2) 从而 (3.1.1), 对射影超曲面 $X$ 成立.
则它对每个有限型分离概形 $X$ 都成立.

\medskip
\noindent\emph{证明.}
由引理~3.8 及 3.6, (3.4.1) 对超曲面的补集 $\mathbb P^n-H$ 成立. 由于

\[
\bigcap_i(\mathbb P^n-H_i)=\mathbb P^n-\bigcup_iH_i
\]

仍是超曲面的补集, $\mathbb P^n-\bigcap_iH_i$ 的开覆盖 $\{\mathbb P^n-H_i\}$
使我们可以应用引理~3.8 及 3.5, 得到 $\mathbb P^n-\bigcap_iH_i$ 的 (3.4.1).
由引理~3.8 及 3.6, (3.4.1) 随后对 $\bigcap_iH_i$ 成立,
因而对每个射影 $X$ 成立. 再次应用引理~3.8 及 3.5,
得到每个拟射影 $X$ 的 (3.4.1), 特别地得到每个仿射 $X$ 的情形.
最后, 若 $X$ 分离且有限型, 取由仿射开子概形组成的有限覆盖并应用引理~3.8 及 3.5 即可.

\medskip
\noindent\emph{底本注.}
印本开头引用了 (3.1.1) 的三个不存在或格式有误的变体;
校订法文层和英文层引用实际公式 (3.1.2)、(3.2.1)--(3.2.2) 及 (3.3.1).
在定理~3.3 的证明中, 印本写作 $\mathrm{fr}_p^{X*}$,
但其旁所列映射为 $\mathrm{fr}_q^X$;
且两次说 Frobenius 在 $q\ne0$ 时将谱序列项化为零,
而应为微分形式次数为正即 $p\ne0$ 时.
最后, 引理~3.9 的证明两次引用不存在的 (3.1.3),
但其论证自身所证明的是 (3.4.1). 照录法文层保留所有这些印本形式;
校订法文层和英文层采用刚才所述的相容读法.
